% TikZ 图：弹簧与胡克定律
% Day 05：力、重力、弹力、摩擦力
% 左：原长，右：悬挂砝码
\begin{tikzpicture}[scale=1,>=Latex]
  % 天花板
  \draw[thick] (-2.5,0) -- (4.5,0);
  
  % 左：原长弹簧
  \draw[decorate,decoration={coil,aspect=0.4,segment length=4pt,amplitude=3pt}] (-1,0) -- (-1,-1.8);
  \node at (-1,-0.9) {$l_0$};
  \node[below] at (-1,-2.2) {原长};
  
  % 右：悬挂砝码
  \draw[decorate,decoration={coil,aspect=0.4,segment length=4pt,amplitude=3pt}] (2.5,0) -- (2.5,-2.8);
  \draw[thick] (2.0,-2.8) rectangle (3.0,-3.3);
  \node at (2.5,-3.05) {$m$};
  
  % 总长标注
  \draw[<->] (3.5,0) -- (3.5,-2.8) node[midway,right,xshift=2pt] {$l$};
  % 伸长量标注（弹性限度范围内额外标出）
  \draw[<->,red] (4.0,-1.8) -- (4.0,-2.8) node[midway,right,xshift=2pt] {$x$};
  
  % 弹力箭头（红，向上）
  \draw[->,thick,red] (2.5,-2.8) -- (2.5,-1.2) node[above right,xshift=2pt] {$F = kx$};
  % 重力箭头（蓝，向下）
  \draw[->,thick,blue] (2.5,-3.3) -- (2.5,-4.2) node[right,xshift=2pt] {$G = mg$};
\end{tikzpicture}
